% ============================================================
% Extended annotated outline (2026-09-27), saved as source material for the
% Rough Draft of Paper (Deliverable 4). NOT the submitted outline.
% To compile, copy into the parent folder temporarily (it uses ../new-aiaa.cls
% and ../ref.bib), since build.ps1 expects one root .tex per folder.
% ============================================================
\documentclass[conf]{new-aiaa}
\usepackage[utf8]{inputenc}

\usepackage{graphicx}
\usepackage{subcaption}
\usepackage{amsmath}
\usepackage[version=4]{mhchem}
\usepackage{siunitx}
\usepackage{float}
\usepackage{array}
\usepackage[export]{adjustbox}
\usepackage{longtable, tabularx}
\usepackage{tabulary}
\usepackage{graphicx,color}
\usepackage{atbegshi,picture}
\usepackage{color, colortbl}
\graphicspath{{./figures/}}
\setlength\LTleft{0pt}

\title{Particle Filter vs. Extended Kalman Filter Localization for a Mobile Manipulator}

\author{Jacob T. Cassady\footnote{Graduate Student, Robotics and Autonomous Systems, AIAA Member.}}
\affil{Johns Hopkins Whiting School of Engineering, Baltimore, Maryland 21218}

\begin{document}

\maketitle
\begin{abstract}
Reliable localization is a critical prerequisite for mobile manipulation, yet algorithms must frequently contend with accumulating odometry drift and perceptual ambiguity in structured environments. This study presents a comparative evaluation of three estimators: dead reckoning, an Extended Kalman Filter (EKF), and a Monte Carlo Localization Particle Filter (PF). These are evaluated on a Clearpath Husky mobile manipulator equipped with a UR5e arm and Robotiq 2F-85 gripper operating in a simulated ROS 2 factory environment. Relying solely on imperfect wheel odometry and 2D LIDAR scans, the filters are evaluated over a three-waypoint route, with 20 seeded trials per configuration, across a one-factor-at-a-time sensitivity sweep of particle count and injected odometry and LIDAR noise. The objective is to quantify the trade-off between the computational efficiency of the EKF during nominal navigation and the robustness of the Particle Filter when recovering from adversarial multimodal scenarios such as global localization and simulated kidnapping.
\end{abstract}

\section{Overview}\label{sec:Overview}
Serving as the paper's introduction, this section will motivate localization as a prerequisite for mobile manipulation: an error in the base pose carries directly into the arm's workspace, and the 0.25~m recovery threshold used throughout will be justified against the UR5e's 0.85~m reach. It will introduce the platform, a Clearpath Husky A200 carrying a UR5e arm and a Robotiq 2F-85 gripper, in a simulated ROS~2/Gazebo warehouse, and explain why localization there is genuinely uncertain inference: the robot's only evidence about its pose is slipping wheel odometry and noisy, ambiguous 2D LIDAR scans. It will then argue why a nonparametric Bayesian filter fits this domain (repetitive geometry, an unknown starting pose, and kidnapping all produce multimodal posteriors that a particle filter can represent and a single-Gaussian EKF cannot) and state the research question as the proposal's two-sided, falsifiable hypothesis: the particle filter recovers in multimodal scenarios where the EKF fails, while the EKF matches its accuracy at lower compute cost during benign runs.

\vspace{0.5em}\noindent\textbf{Planned Literature Integration:}
\begin{itemize}
    \item \textbf{Macenski, S., et al. (2022) \cite{macenski2022ros2}:} This paper will be used for \textbf{background} on the ROS~2 node, topic, and bridge architecture that connects the simulator, the sensors, and the estimators.
    \item \textbf{Koenig, N., \& Howard, A. (2004) \cite{koenig2004design}:} This paper will be used for \textbf{background} to introduce Gazebo, the simulator whose physics produces the wheel slip in the odometry and whose pose publisher supplies the ground truth used for scoring.
\end{itemize}

\section{Problem Statement}\label{sec:Problem-Statement}
This section will formalize localization as recursive Bayesian estimation of the planar pose $(x, y, \theta)$ from odometry and LIDAR scans against a known occupancy map, and will name three specific sources of uncertainty. First, \textbf{odometry drift}: the drive controller integrates wheel-encoder rates on a skid-steer base, so wheel slip generated by Gazebo's wheel-slip system, together with the injected odometry noise defined in Sec.~\ref{sec:Protocol}, makes dead-reckoned pose error grow without bound. Second, \textbf{perceptual ambiguity}: Gazebo's LIDARs return noise-free ranges, so range noise is injected at the levels in Sec.~\ref{sec:Protocol}, and the warehouse's repeated identical shelving units can make one scan consistent with several distinct poses. Third, \textbf{kidnapping}: an unmodeled displacement of the robot (by Gazebo teleportation) that, like the unknown initial pose of global localization, leaves the true pose far from the current estimate and makes the posterior multimodal, violating the single-Gaussian assumption; these two cases are exactly the adversarial tests of Sec.~\ref{sec:Protocol}.

\vspace{0.5em}\noindent\textbf{Planned Literature Integration:}
\begin{itemize}
    \item \textbf{Thrun, S., Burgard, W., \& Fox, D. (2005) \cite{thrun2005probabilistic}:} This text will be used for \textbf{background}, providing the taxonomy of localization problems (position tracking, global localization, and the kidnapped robot problem) that organizes the three sources of uncertainty and motivates the adversarial tests.
    \item \textbf{ssarkar (2024) \cite{ssarkar2024warehouse}:} This Gazebo Fuel world, assembled from AWS RoboMaker warehouse models, will be used as the \textbf{data} source: it is the environment in which every trial is recorded, and its repeated shelving is the origin of the perceptual ambiguity studied.
\end{itemize}

\section{Study Approach}\label{sec:Study-Approach}
Detailing the methodology, this section will define the platform and data pipeline, the three estimators, how the comparison is kept fair despite their different origins, and the experimental protocol and statistics. Each subsection below corresponds to one of these parts.

\subsection{Simulation Platform and Data Collection}\label{sec:Platform}
This subsection will describe the \texttt{devol} model: a Clearpath Husky A200 with a UR5e arm and Robotiq 2F-85 gripper, carrying two planar LIDARs, a front-mounted $270^\circ$ scanner (540 beams, 0.05--25~m, 40~Hz) and an elevated $360^\circ$ scanner on the sensor arch (900 beams, 0.9--130~m, 20~Hz). The shared map will be generated with \texttt{octomap\_server} from a point cloud of the warehouse's static geometry and projected to a 2D occupancy grid, so that map error does not confound the comparison. Gazebo's pose publisher supplies ground truth at 50~Hz, which drives the waypoint controller so estimator error never alters the trajectory being scored. Because the estimators are therefore passive observers, each trial's odometry, scans, and ground truth are recorded once and replayed offline through every estimator and configuration, which guarantees identical inputs.

\vspace{0.5em}\noindent\textbf{Planned Literature Integration:}
\begin{itemize}
    \item \textbf{Hornung, A., et al. (2013) \cite{hornung2013octomap}:} This paper will be used for \textbf{methodology}, describing the probabilistic octree representation that \texttt{octomap\_server} builds from the static point cloud and from which the 2D occupancy grid is projected.
\end{itemize}

\subsection{Baseline Estimators: Dead Reckoning and EKF}\label{sec:Baselines}
Dead reckoning integrates the recorded odometry from the known start pose and serves as the lower bound that any LIDAR correction must beat. The EKF, adapted from Introduction to Robotics classwork, tracks $(x, y, \theta)$ with an odometry-driven prediction, whose process noise is derived from the same motion-noise parameters $\alpha_{1..4}$ as the particle filter, and an update from a pose observation produced by matching each scan against the map with an ICP-based scan matcher, linearizing both models about the current mean. For global localization, the EKF is initialized at the centroid of the free space with a covariance spanning the map, so its expected failure reflects the single-Gaussian representation rather than an arbitrary initial guess.

\vspace{0.5em}\noindent\textbf{Planned Literature Integration:}
\begin{itemize}
    \item \textbf{Besl, P. J., \& McKay, N. D. (1992) \cite{besl1992method}:} This paper will be used for \textbf{methodology}, describing the ICP registration algorithm underlying the scan-to-map matcher that produces the EKF's pose observation.
    \item \textbf{Moore, T., \& Stouch, D. (2014) \cite{moore2014generalized}:} This paper describes the generalized EKF in the ROS \texttt{robot\_localization} package and will be used for \textbf{background} on EKF-based localization in ROS, situating the adapted classwork EKF among standard implementations.
\end{itemize}

\subsection{Particle Filter (Monte Carlo Localization)}\label{sec:PF}
The particle filter implements the bootstrap (SIR) form of Monte Carlo Localization: each of $N$ particles is propagated by sampling the odometry motion model, weighted by a likelihood-field measurement model (a mixture of hit, random, and max-range components) evaluated against the map, and resampled with the low-variance sampler. Because a plain SIR filter cannot recover once its particles have collapsed around a wrong pose, the implementation adds the Augmented MCL recovery step, which injects uniformly drawn particles in proportion to how far the short-term average measurement likelihood has fallen below its long-term average.

\vspace{0.5em}\noindent\textbf{Planned Literature Integration:}
\begin{itemize}
    \item \textbf{Gordon, N. J., Salmond, D. J., \& Smith, A. F. M. (1993) \cite{gordon1993novel}:} This paper introduced the bootstrap filter and will be used for \textbf{methodology}, defining the propagate--weight--resample cycle the particle filter implements.
    \item \textbf{Dellaert, F., Fox, D., Burgard, W., \& Thrun, S. (1999) \cite{dellaert1999monte}:} This paper introduced Monte Carlo Localization and will be used for \textbf{methodology}, as the formulation that applies the bootstrap filter to pose estimation against a known map.
    \item \textbf{Thrun, S., Burgard, W., \& Fox, D. (2005) \cite{thrun2005probabilistic}:} This text will be used for \textbf{methodology}, providing the odometry motion model, the likelihood-field measurement model, the low-variance sampler, the Augmented MCL recovery step, and the EKF localization derivation used in Sec.~\ref{sec:Baselines}.
    \item \textbf{Arulampalam, M. S., et al. (2002) \cite{arulampalam2002tutorial}:} This tutorial will be used for \textbf{background} on sequential Monte Carlo methods, justifying the selection of the bootstrap (SIR) form and its resampling step.
    \item \textbf{Lenser, S., \& Veloso, M. (2000) \cite{lenser2000sensor}:} This paper's sensor-resetting localization will be used for \textbf{comparison}, as the alternative recovery mechanism against which the choice of Augmented MCL's random-particle injection is justified.
\end{itemize}

\subsection{Implementation Asymmetry}\label{sec:Asymmetry}
This subsection will acknowledge that the particle filter is written from scratch for this study while the EKF is adapted from prior classwork, and that the two differ in how they use the scan: the particle filter scores every beam against the likelihood field, whereas the EKF compresses each scan into a single pose observation from the scan matcher. To keep the comparison about the posterior representation rather than the code, both filters receive identical recorded inputs, the same map, and motion noise derived from the same $\alpha$ values. Because of this difference, an EKF failure could stem either from the scan matcher converging to a wrong pose or from the single-Gaussian assumption itself; Sec.~\ref{sec:Limitations} weighs these explanations against the results.

\subsection{Experimental Protocol and Sensitivity Sweep}\label{sec:Protocol}
Each estimator is scored on the three-waypoint factory route over 20 trials per configuration, where each trial's seed sets the injected-noise realization and the particle filter's sampling; because Gazebo's sensors are noise-free and its physics nearly deterministic, the seed, not the simulator, is the source of trial-to-trial variation. Noise is injected offline into the recorded streams so each level is exactly controlled: odometry increments are perturbed with the odometry motion model using $\alpha_1 = \alpha_2 = \alpha_3 = \alpha_4 = k\,\alpha_{\text{nom}}$ with $\alpha_{\text{nom}} = 0.05$ and $k \in \{0, 1, 2, 4\}$ ($k = 0$ leaves only the physics-induced wheel slip), and every LIDAR range receives zero-mean Gaussian noise with $\sigma_r \in \{0.01, 0.03, 0.10\}$~m, while the filters' own noise parameters stay at their nominal values so that higher levels test robustness to under-modeled noise. Particle count $N \in \{100, 500, 2000, 5000\}$, $k$, and $\sigma_r$ are swept one factor at a time about the nominal point ($N = 2000$, $k = 1$, $\sigma_r = 0.03$~m), giving nine particle-filter, six EKF, and four dead-reckoning configurations. The adversarial tests run 20 trials each at the nominal point: (a) global localization, with the particle filter initialized from a uniform prior over free space, and (b) kidnapping, in which the robot is teleported mid-route to a different aisle without the estimators being told.

\subsection{Metrics and Statistical Analysis}\label{sec:Statistics}
Metrics are position RMSE, heading RMSE (with angle wrapping), position error at each of the three waypoints, and wall-clock compute time per filter update on a single fixed machine. Twenty trials per configuration bound each reported mean to a 95\% confidence-interval half-width of $t_{0.975,19}\,s/\sqrt{20} \approx 0.47\,s$, where $s$ is the between-trial standard deviation. For the binary recovery outcome, 20/20 successes gives a Clopper--Pearson 95\% lower bound of 0.83 on the recovery probability and 0/20 an upper bound of 0.17, so twenty trials separate a filter that reliably recovers from one that reliably does not.

\vspace{0.5em}\noindent\textbf{Planned Literature Integration:}
\begin{itemize}
    \item \textbf{Fox, D. (2001) \cite{fox2001kld}:} This paper's KLD-sampling, which adapts the particle count online, will be used for \textbf{comparison}, framing the fixed-$N$ sweep and the accuracy--compute trade-off it measures.
\end{itemize}

\section{Results}\label{sec:Results}
This section will report every metric as mean $\pm$ 95\% confidence interval over the 20 trials, organized to answer the two halves of the hypothesis in turn.

\subsection{Nominal-Route Accuracy and Compute}\label{sec:Results-Nominal}
A table will report position RMSE, heading RMSE, waypoint error, and per-update compute time for all three estimators at the nominal point. A companion figure will plot position error against time for one representative trial, contrasting dead reckoning's unbounded drift with the two filters' bounded error. This subsection answers the benign half of the hypothesis: whether the EKF matches or exceeds the particle filter's accuracy at lower compute cost.

\subsection{Sensitivity to Particle Count and Injected Noise}\label{sec:Results-Sensitivity}
Figures will plot particle-filter RMSE and per-update compute time against $N$ on a logarithmic axis, with confidence bands and the EKF as a horizontal reference, locating the particle count at which the particle filter matches the EKF's nominal accuracy. Further figures will plot RMSE against $k$ and $\sigma_r$ for all three estimators, showing how gracefully each degrades as the evidence becomes noisier than its model assumes.

\vspace{0.5em}\noindent\textbf{Planned Literature Integration:}
\begin{itemize}
    \item \textbf{Fox, D., Burgard, W., Dellaert, F., \& Thrun, S. (1999) \cite{fox1999monte}:} This paper will be used for \textbf{comparison}: its reported relationship between sample-set size and localization accuracy, and its global-localization results, are the published reference against which the $N$ sweep and the recovery behavior are compared.
\end{itemize}

\subsection{Adversarial Recovery}\label{sec:Results-Adversarial}
For global localization and kidnapping, this subsection will report each estimator's recovery rate with Clopper--Pearson intervals and its recovery time, defined as the simulated time until position error falls below 0.25~m and remains below it. Trials that do not recover within a fixed timeout count as failures, and recovery time is summarized over successful trials only, so the two numbers are read together. A figure will overlay position error against time around the kidnapping event for all three estimators, answering the adversarial half of the hypothesis.

\subsection{Limitations and Alternative Explanations}\label{sec:Limitations}
This subsection will connect specific results to the study's limitations. For any EKF failure after kidnapping, it will weigh the counter-explanation raised in Sec.~\ref{sec:Asymmetry}: that the scan matcher, not the Gaussian assumption, caused it. It will also note that injected Gaussian range noise is kinder than real LIDAR returns, which include glass, multipath, and dynamic clutter, and that the ambiguity findings are specific to a single warehouse layout.

\section{Conclusion}\label{sec:Conclusion}
The final section will state whether the evidence supports each half of the hypothesis, citing the specific recovery rates and the nominal-route accuracy and compute results that decide it. It will discuss the practical implications for the platform's waypoint controller and planning pipeline, such as whether the particle filter's extra compute is justified only for relocalization, which would suggest a hybrid in which the EKF tracks and the particle filter re-initializes it. It will state what would need to change for the conclusions to generalize beyond this simulation: real sensor noise, dynamic obstacles, and other map geometries. Finally, it will propose future work, including an EKF and particle filter that share the same measurement front end to remove the implementation asymmetry, and KLD-sampling~\cite{fox2001kld} to adapt $N$ online.

\section*{Acknowledgments}\label{sec:Acknowledgments}
Anthropic's Claude was used to review this outline against the course rubric and the instructor's feedback on the proposal, and to draft revisions to it. The final paper will extend this disclosure to any AI assistance with its code or writing.

\bibliography{ref}

\end{document}
